\documentclass[10pt, a4paper]{article}
%-----------------------
%- 	PACKAGES & SETTINGS
%-----------------------
\usepackage[a4paper,top=1.5cm,bottom=2cm,left=4cm,right=4cm]{geometry}
\usepackage[utf8]{inputenc}
\usepackage[italian]{babel}
\usepackage{xcolor}
\usepackage{hyperref}
\hypersetup{
    colorlinks=true,
    filecolor=magenta,      
    urlcolor=darkgray,
    linkcolor=black
}
\urlstyle{same}
\usepackage{amsmath}
\usepackage{graphicx}
\graphicspath{ {images/} }
 
%-----------------------
%- 	TITLE
%-----------------------
\title{\textbf{Relazione Progetto OCaml PR 2}}
\author{\textbf{Venturi} Ludovico\\Docente: \href{http://pages.di.unipi.it/levi/}{Francesca Levi}}
\date{UNIPI, Gennaio 2020}


%-----------------------
%- 	DOCUMENT
%-----------------------
\begin{document}
%- 	INTRO
\pagenumbering{roman} 
\maketitle
\tableofcontents
\vfill
\begin{figure}[h]
	\centering
	\includegraphics[scale=0.3]{ocaml_logo}
	\label{fig:0}
\end{figure}

\clearpage

%- 	START DOC
\pagenumbering{arabic} 
\section{Scelte progettuali}
\begin{itemize}
\item Tutti i valori associati alle chiavi in un dizionario hanno lo stesso tipo
\item 
\end{itemize}
\subsection{Sintassi Astratta}
Estensione della sintassi astratta del linguaggio  \texttt{type exp = }\\
\texttt{...| Dict of (ide * exp)list 			| Insert of ide * exp * exp |
	Delete of  ide * exp 				| HasKey of ide * exp				| Iterate of exp * exp	|
	Fold of exp * exp					| Filter of ide list * exp  		| PrintDictValues of exp;;}\\\\
Estensione dei tipi esprimibili	\texttt{type evT = }\\
\texttt{...|DictVal of (ide * evT)list}
\section{Utilizzo}
Aprire l'interprete top-level di OCaml da terminale digitando \texttt{ocaml}. Importare l'interprete del nostro linguaggio:
\begin{center}
\texttt{\# use "venturi.ml";;}
\end{center}
e da qui valutare le espressioni, entrando nel vivo del \textbf{REPL}(ReadEvalPrintLoop).
\subsection{Esempi}
Assumendo \texttt{let env0 : ide -> evT = emptyenv Unbound;;}
\begin{itemize}
\item \texttt{\# let x = Dict([("Birman", Eint(3));("Mainecoon", Eint(13));("Siamese", Eint(17));("Foldex", Eint(21))]);;}\\
\textit{val x : exp =
  Dict
   [("Birman", Eint 3); ("Mainecoon", Eint 13); ("Siamese", Eint 17);
    ("Foldex", Eint 21)]
}
\item \texttt{\# let y = eval (Insert("Korat", Eint(4), x)) env0;; }\\
\textit{val y : evT =
  DictVal
   [("Korat", Int 4); ("Birman", Int 3); ("Mainecoon", Int 13);
    ("Siamese", Int 17); ("Foldex", Int 21)]
}
\end{itemize}
\end{document}









